\subsection{Annual estimands and sample supports}
\label{sec:si-annual-models}

The annual analysis distinguishes recorded activity from facility exposure.  For
city $i$ and year $t$, let $Q_{it}$ be recorded charging activity, $E_{it}$
the cumulative facility indicator summed over the same compatible weeks, and
$W_{it}$ the number of compatible weeks.  The source-defined measures are
\begin{equation}
 A_{it}=Q_{it}/W_{it},\qquad N_{it}=E_{it}/W_{it},\qquad
 I_{it}=Q_{it}/E_{it}=A_{it}/N_{it},\qquad x_{it}=\log N_{it}.
 \label{eq:si-definitions}
\end{equation}
Compatibility requires a valid source date and positive facilities.  A
source-coded zero activity is retained, whereas a city-year with no compatible
weeks is missing and is not assigned a zero.  Thus $A$ is activity per
compatible week, $N$ is time-average facility exposure, and $I$ is activity per
cumulative facility exposure.  These ratios do not recover latent demand,
instantaneous online capacity or a complete operator total.

The annual sample and its exclusions are shown in Table~\ref{tab:sample_models}.
Before fixed-effect support exclusions, the positive-support sample contains
464 $Q>0$, $E>0$, $W>0$ rows from 123 cities over 2016--2023.  Nineteen
cities contribute exactly one positive year; those rows provide no within-city
information for a city fixed-effect slope and are excluded from the positive
models, leaving 445 rows in 104 cities.  The zero-inclusive support begins
with 887 rows from 292 cities with finite nonnegative $Q$, positive $E$ and
$W$, positive finite $N$, and complete annual weather controls.  It excludes
169 cities (353 rows) whose support-wide activity sum is zero and then 11
singleton cities (11 rows), leaving 523 rows from 112 cities: 453 positive
rows and 70 source-coded-zero rows.  These are support definitions, not
imputations of an unobserved outcome.

\subsection{Conditional means and model specifications}
For the log outcome, the constant-elasticity specification is
\begin{equation}
 \operatorname{E}[\log A_{it}\mid X_{it}]
 =\alpha_i+\lambda_t+\beta x_{it}+\gamma^{\top}Z_{it},
 \label{eq:si-log-l0}
\end{equation}
and a restricted-cubic extension adds two nonlinear terms:
\begin{equation}
 \operatorname{E}[\log A_{it}\mid X_{it}]
 =\alpha_i+\lambda_t+\beta x_{it}+
 \theta_1H_1(x_{it})+\theta_2H_2(x_{it})+\gamma^{\top}Z_{it}.
 \label{eq:si-log-models}
\end{equation}
Here $\alpha_i$ and $\lambda_t$ are city and calendar-year fixed effects.  The
vector $Z_{it}$ contains annual HDD18, CDD24, precipitation, mean wind speed
and mean relative humidity.  These aggregates adjust for measured annual
background conditions; they do not make facility growth exogenous or provide an
instrument.

The level models have a different estimand and retain observation exposure:
\begin{equation}
 \operatorname{E}[Q_{it}\mid X_{it},W_{it}]
 =W_{it}\exp\!\left(\alpha_i+\lambda_t+\beta x_{it}
 +\gamma^{\top}Z_{it}\right),
 \label{eq:si-p0}
\end{equation}
and
\begin{equation}
 \operatorname{E}[Q_{it}\mid X_{it},W_{it}]
 =W_{it}\exp\!\left(\alpha_i+\lambda_t+\beta x_{it}
 +\theta_1H_1(x_{it})+\theta_2H_2(x_{it})
 +\gamma^{\top}Z_{it}\right).
 \label{eq:si-p1}
\end{equation}
The two level equations use Poisson pseudo-maximum likelihood (PPML), whose
estimand is the conditional arithmetic mean of nonnegative $Q$; integer
outcomes and equidispersion are not required.  The positive-support level
equations use the common 445-row positive support, whereas the zero-inclusive
level equation uses the separate 523-row support.  A zero outcome in the zero-inclusive sample is a
source-coded zero under the stated compatibility and weather rules; unknown
activity is never treated as zero.  Panels (a) and (b) of
Figure~\ref{fig:si_annual} use the 445-row, 104-city positive-support fits;
panel (c) additionally compares the positive-support and zero-inclusive
constant PPML estimates.  The two models therefore do not estimate a single
relation in which missingness and a recorded zero are interchangeable.

The offset is $\log W_{it}$ in the PPML linear predictor, not an additional
frequency-weight argument.  If $\mu_{it}$ denotes the modeled weekly mean and
$X_{it}$ is a design row, then
\begin{equation}
 \operatorname{E}[Q_{it}\mid X_{it},W_{it}]=W_{it}\mu_{it},\qquad
 X_{it}(Q_{it}-W_{it}\mu_{it})
 =W_{it}X_{it}(A_{it}-\mu_{it}).
 \label{eq:si-offset-equivalence}
\end{equation}
Compatible-week exposure therefore changes the modeled mean and the
score/information; the absence of a frequency weight does not imply equal
regression influence.  The PPML coefficient is interpreted on the conditional
arithmetic-mean scale, not as an OLS coefficient for a log-transformed
outcome.

\subsection{Restricted-cubic form and local elasticity}
The log and level specifications use the same two-term restricted-cubic basis.
Let $k_1<k_2<k_3<k_4$ be the empirical 0.05, 0.35, 0.65 and 0.95 quantiles of
$x$ in the 445-row positive sample.  Linear quantile interpolation gives the following display-rounded values
\[
(k_1,k_2,k_3,k_4)=(1.0986,\ 3.3273,\ 4.7236,\ 7.3771).
\]
With $(u)_+=\max(u,0)$,
\begin{equation}
 H_j(x)=\frac{(x-k_j)_+^3-
 \dfrac{k_4-k_j}{k_4-k_3}(x-k_3)_+^3+
 \dfrac{k_3-k_j}{k_4-k_3}(x-k_4)_+^3}
 {(k_4-k_1)^2},\qquad j\in\{1,2\}.
 \label{eq:si-rcs}
\end{equation}
The analytic derivative used for local elasticity is
\begin{equation}
 H'_j(x)=\frac{3(x-k_j)_+^2-
 3\dfrac{k_4-k_j}{k_4-k_3}(x-k_3)_+^2+
 3\dfrac{k_3-k_j}{k_4-k_3}(x-k_4)_+^2}
 {(k_4-k_1)^2}.
 \label{eq:si-rcs-derivative}
\end{equation}
For either fitted log mean, the local elasticity is
$\eta(x)=\beta+\theta_1H'_1(x)+\theta_2H'_2(x)$.  It is the local derivative
of the stated conditional function, not an estimated threshold or optimal
stock.  Curve intervals are pointwise city-cluster intervals; no simultaneous
band is claimed.  The joint restrictions remain model-specific: for the
nonlinear log specification, Wald $=2.701$, $F(2,103)=1.351$ and
$p=\LogSplineP$; for the nonlinear level specification, Wald $=3.021$,
$F(2,103)=1.510$ and $p=\MeanSplineP$.  A non-rejection concerns the two
added basis terms on that model's own outcome scale; it is not a global
linearity test or an equality test across models.

\subsection{Clustered inference and contrasts}
For OLS, city-clustered CR1 inference uses the full design rank:
\begin{equation}
 \widehat V_{\mathrm{CR1}}=
 \frac{G}{G-1}\frac{n-1}{n-r}
 (X^\top X)^{-1}
 \left[\sum_{g=1}^{G}X_g^\top\hat u_g\hat u_g^\top X_g\right]
 (X^\top X)^{-1}.
 \label{eq:si-cr1-ols}
\end{equation}
Here $r$ is the actual matrix rank.  The PPML sandwich uses the Poisson
Hessian $X^\top\operatorname{diag}(\hat\mu)X$ and city-summed scores
$X_i(Q_i-\hat\mu_i)$.  The positive sample has rank 117 for the linear log and
positive-support level models and rank 119 for their nonlinear extensions, with
$n=445$, $G=104$, CR1 factors 1.3668 and 1.3752, and $t(G-1)=t(103)$
intervals.  The zero-inclusive level model has rank 125, $G=112$ and
$t(111)$ intervals.  For a local elasticity, if $d(x)$ contains
$1,H'_1(x),H'_2(x)$ in the relevant scale-coefficient positions, its variance
is $d(x)'\widehat Vd(x)$ and its pointwise interval is
\[
 \hat\eta(x)\mathbin{\pm}t_{G-1,.975}\sqrt{d(x)'\widehat Vd(x)}.
\]
The joint test uses $F=(\theta'V_\theta^{-1}\theta)/2$ with its
model-specific covariance block.  These conditional intervals address
within-city dependence, not sample selection or policy endogeneity.
The main annual table (Table~\mainref{tab:annual_models}) reports the three
constant-elasticity estimates, their $\beta-1$ contrasts and the corresponding
tests against zero and one.  The rank, cluster degrees of freedom and
covariance construction above supply the inference details for those reported
contrasts, with each test remaining on its model's stated conditional-mean scale.
